% =============================================================================
% Retrocopertina
% =============================================================================
\pagestyle{empty}
\cleardoublepage
\begin{tikzpicture}[remember picture,overlay]
	\fill[modernblack] (current page.south west) rectangle (current page.north east);
	\node[anchor=north, text width=0.8\textwidth, align=center]
	at ([yshift=-2cm] current page.north) {%
		\color{white}
		{\sffamily\Large\textit{``La tesi del libro, in una frase.''}}\\[0.8em]
		{\sffamily\normalsize\textls[50]{\MakeUppercase{\authorname}}}
	};
	\node[anchor=center, text width=0.8\textwidth, align=justify]
	at (current page.center) {%
		\sffamily\small\color{white}
		Il problema, la tesi e i tre argomenti principali in poche righe:
		la quarta di copertina è un piccolo testo argomentativo.
	};
\end{tikzpicture}
